\documentclass[11pt]{article}
\usepackage[margin=1in]{geometry}
\usepackage{amsmath,amssymb,amsthm}
\usepackage{hyperref}
\newtheorem{theorem}{Theorem}
\newtheorem{lemma}[theorem]{Lemma}
\newtheorem{corollary}[theorem]{Corollary}
\newtheorem{proposition}[theorem]{Proposition}
\theoremstyle{remark}
\newtheorem{remark}[theorem]{Remark}
\newcommand{\F}{\mathbb{F}_3}
\newcommand{\wt}{\operatorname{wt}}

\title{Disproof of Conjecture 00000001090:\\
there is no MDS code of type $[12,6]$ over $\mathrm{GF}(3)$}
\author{jilint777}
\date{2026-10-04}

\begin{document}
\sloppy
\maketitle

\begin{abstract}
Conjecture 00000001090 asserts that the maximal order of the automorphism group
of an MDS code over $\mathrm{GF}(3)$ of type $[12,6]$ is $660$, attained by an
$M_{11}$ embedding, and that the orders for all other equivalence classes divide
$132$. We prove that \emph{no} ternary MDS code of type $[12,6]$ exists. More
generally, if a $k\times n$ matrix over $\F$ with $k\ge2$ has all nonzero
combinations of its rows of weight at least $n-k+1$, then $n\le k+2$. Hence
no code attains the order $660$, and the conjecture is false whatever notion of
automorphism group is used. If the word ``MDS'' is dropped, the claim is still
false: the ternary $[12,6]$ code of repeated pairs has $2^6\cdot6!=46080>660$
permutation automorphisms. The main results are formalized in Lean~4 (core
library only).
\end{abstract}

\section{The conjecture and how we read it}
\begin{quote}
\emph{Definition:} The automorphism group of an MDS code over $\mathrm{GF}(3)$
of type $[12,6]$ is the stabilizer of its permutation equivalence class.
\emph{Conjecture:} The maximal order of this automorphism group is $660$
(attained by an $M_{11}$ embedding); the group orders of all other equivalence
classes divide $132$.
\end{quote}
The Chinese version says the same.

A ternary linear code of type $[n,k]$ is a $k$-dimensional subspace
$C\subseteq\F^n$. Its minimum distance $d$ is the least Hamming weight of a
nonzero codeword. The Singleton bound gives $d\le n-k+1$, and $C$ is
\emph{MDS} (maximum distance separable) if $d=n-k+1$. So an MDS code of type
$[12,6]$ is a ternary $[12,6,7]$ code.

The conjecture is a conjunction. Its first clause says that the maximum of
$|\mathrm{Aut}(C)|$ over MDS $[12,6]$ ternary codes $C$ \emph{is} $660$, and that
this value is \emph{attained}. In particular it asserts that there is an MDS
$[12,6]$ ternary code whose automorphism group has order $660$. We refute this
clause, which refutes the conjecture. The remaining clauses are about all MDS
$[12,6]$ ternary codes. Since there are none, they are vacuously true, and we
make no claim about them.

The definition of the automorphism group is worded unusually. The usual
meaning is the stabilizer of $C$ in the group acting on codes: permutations of
coordinates (the permutation group $\mathrm{PAut}(C)$), or monomial maps
(permutations combined with nonzero scalars, giving $\mathrm{MAut}(C)$). Over
$\mathrm{GF}(3)$ there are no nontrivial field automorphisms, so $\mathrm{MAut}$
is also the full automorphism group. Our main result is that the set of MDS
$[12,6]$ ternary codes is empty, so the refutation does not depend on which
notion is meant (Section~\ref{sec:readings}).

\section{Notation}
Write $\F=\{0,1,2\}$ with arithmetic mod $3$. Every nonzero $x\in\F$ satisfies
$x^2=1$. For $w\in\F^n$, $\wt(w)$ is the number of nonzero coordinates. A
$k\times n$ matrix $G$ with rows $g_0,\dots,g_{k-1}$ sends a message
$m\in\F^k$ to the codeword $mG=\sum_i m_ig_i$. The rows are linearly
independent if $mG\ne0$ for all $m\ne0$. In that case $C=\{mG\}$ is an
$[n,k]$ code. Every $[n,k]$ code arises this way (take a basis as the rows).

\section{The main theorem}
\begin{lemma}[linear dependence]\label{lem:dep}
Let $r<k$ and let $v_0,\dots,v_{k-1}\in\F^r$. There is a nonzero
$c\in\F^k$ with $\sum_i c_iv_i=0$.
\end{lemma}
\begin{proof}
Induction on $r$. For $r=0$ take $c=(1,0,\dots,0)$. Let $r\ge1$ and look at the
last vector $w=v_{k-1}$. If $w=0$, take $c=(0,\dots,0,1)$. Otherwise choose a
coordinate $j_0$ with $p=w_{j_0}\ne0$, and put
$v_i'=v_i-(v_{i,j_0}p)\,w$ for $i<k-1$. Since $p^2=1$, every $v_i'$ vanishes at
$j_0$. Deleting coordinate $j_0$ gives $k-1>r-1$ vectors of $\F^{r-1}$. By
induction there is a nonzero $c'\in\F^{k-1}$ with $\sum_ic'_iv'_i=0$. Then
$c=(c'_0,\dots,c'_{k-2},\,-\sum_ic'_iv_{i,j_0}p)$ is nonzero and
$\sum_ic_iv_i=\sum_ic'_iv'_i=0$.
\end{proof}

\begin{lemma}[counting]\label{lem:count}
For $x,y\in\F$ let $N(x,y)$ be the number of the eight pairs
$(a,b)\in\F^2\setminus\{(0,0)\}$ with $ax+by\ne0$. Then $N(0,0)=0$ and
$N(x,y)=6$ otherwise.
\end{lemma}
\begin{proof}
If $(x,y)\ne(0,0)$, the map $(a,b)\mapsto ax+by$ is a nonzero linear form on
$\F^2$. Its kernel has $3$ elements, one of which is $(0,0)$. So $9-3=6$ of the
eight nonzero pairs give a nonzero value.
\end{proof}

\begin{theorem}\label{thm:main}
Let $k\ge2$ and let $G$ be a $k\times n$ matrix over $\F$ such that every
nonzero message $m$ gives $\wt(mG)\ge n-k+1$. Then $n\le k+2$.
\end{theorem}
\begin{proof}
Suppose $n\ge k+3$.

\emph{Step 1.} Apply Lemma~\ref{lem:dep} to the restrictions of the rows
$g_0,\dots,g_{k-1}$ to the first $k-2$ coordinates. This gives a nonzero
message $m^{(1)}$ such that $u=m^{(1)}G$ vanishes on coordinates
$0,\dots,k-3$. Fix $i_0$ with $m^{(1)}_{i_0}\ne0$.

\emph{Step 2.} Apply Lemma~\ref{lem:dep} with $r=k-1$ to the vectors
$(g_{i,0},\dots,g_{i,k-3},\,[i=i_0])\in\F^{k-1}$. Here the last entry is $1$ for
$i=i_0$ and $0$ otherwise. This gives a nonzero message $m^{(2)}$ such that
$v=m^{(2)}G$ vanishes on coordinates $0,\dots,k-3$, and $m^{(2)}_{i_0}=0$.

\emph{Step 3.} Let $(a,b)\ne(0,0)$. The message $m=am^{(1)}+bm^{(2)}$ is
nonzero: if $a\ne0$ then $m_{i_0}=am^{(1)}_{i_0}\ne0$; if $a=0$ then $b\ne0$
and $m=bm^{(2)}\ne0$. Hence $\wt(au+bv)=\wt(mG)\ge n-k+1$. Summing over the
eight nonzero pairs and exchanging the sums,
\[
8(n-k+1)\le\sum_{(a,b)\ne(0,0)}\wt(au+bv)=\sum_{j=0}^{n-1}N(u_j,v_j)
\le 6\,(n-k+2),
\]
by Lemma~\ref{lem:count}, because $u_j=v_j=0$ for the first $k-2$ coordinates.
So $2(n-k)\le4$, that is $n\le k+2$, a contradiction.
\end{proof}

Theorem~\ref{thm:main} is the case $q=3$ of the classical bound $n\le q+k-1$
for MDS codes with $k\ge2$ \cite{ms}. In coding-theoretic language, Steps 1--2
shorten the code $k-2$ times, producing a two-dimensional code of length
$n-k+2$ with minimum distance at least $n-k+1$. The counting of Step~3 is
the Plotkin argument for that two-dimensional code. The bound is sharp: the
tetracode generated by $(1,0,1,1)$ and $(0,1,1,2)$ is a $[4,2,3]$ MDS code.

\begin{corollary}\label{cor:main}
There is no ternary $[12,6,7]$ code. More strongly, no $6\times12$ matrix over
$\F$ has all nonzero combinations of its rows of weight $\ge7$. Likewise, there
is no ternary $[12,7,6]$ code.
\end{corollary}
\begin{proof}
Take $(n,k)=(12,6)$, respectively $(12,7)$, in Theorem~\ref{thm:main}:
$12>k+2$.
\end{proof}

Concretely, for $[12,6]$: there are $9$ codewords vanishing on the first four
coordinates, forming a $2$-dimensional space. Restricted to the last $8$
coordinates, the $8$ nonzero ones would all need weight $\ge7$, total $\ge56$. But
each of the $8$ coordinates contributes at most $6$, so the total is at most $48$.
The Griesmer bound gives another proof: a ternary $[n,6,7]$ code needs
$n\ge\sum_{i=0}^{5}\lceil7/3^i\rceil=14$.

\section{Consequences for the conjecture}\label{sec:readings}
\begin{theorem}
Conjecture 00000001090 is false.
\end{theorem}
\begin{proof}
By Corollary~\ref{cor:main} there is no MDS code over $\mathrm{GF}(3)$ of type
$[12,6]$. So no such code has an automorphism group of order $660$, and the
first clause (``the maximal order is $660$, attained'') fails.
\end{proof}

This argument covers every reasonable reading:
\begin{itemize}
\item \emph{Notion of automorphism group.} Permutation group, monomial group,
full automorphism group, or literally ``the stabilizer of the permutation
equivalence class'': none of them matter, because there is no code. In Lean we
state the conjecture for $\mathrm{PAut}$ and for $\mathrm{MAut}$. We also prove
that no MDS $[12,6]$ ternary code has \emph{any} property $P$.
\item \emph{Definition of MDS.} We only use $d\ge n-k+1=7$, so the strict
($d=7$) and non-strict readings are both covered.
\item \emph{Type $[12,6]$ as $[n,d]$.} If $[12,6]$ meant length $12$ and minimum
distance $6$, an MDS code would be $[12,7,6]$, which does not exist either
(Corollary~\ref{cor:main}).
\item \emph{``MDS'' dropped.} Suppose the claim is read for all ternary
$[12,6]$ codes, for example because the author had the ternary Golay code in
mind. It is still false. The code of repeated pairs
\[
P=\{(x_0,x_0,x_1,x_1,\dots,x_5,x_5)\}
\]
is a $[12,6,2]$ code. Its weight-$2$ words have exactly the supports
$\{2i,2i+1\}$, so every automorphism permutes these six pairs. Conversely,
every permutation preserving the pair partition is an automorphism. So
$|\mathrm{PAut}(P)|=2^6\cdot6!=46080>660$. Lean exhibits $768$ distinct
permutation automorphisms of $P$. For the monomial group the order is even larger,
since $\mathrm{PAut}\subseteq\mathrm{MAut}$. (The ternary Golay code
$[12,6,6]$ is also well known to have monomial automorphism group $2.M_{12}$ of
order $190080$ \cite{cs}. We cite this but do not verify it.)
\end{itemize}

\begin{remark}
The statement is also internally inconsistent: $|M_{11}|=7920$, so $M_{11}$
does not embed in a group of order $660$. ($660=|\mathrm{PSL}(2,11)|$, and
$\mathrm{PSL}(2,11)$ is a maximal subgroup of $M_{11}$.) We do not use this.
\end{remark}

\begin{remark}[non-vacuity]
The obstruction is exactly the distance $7$. Ternary $[12,6]$ codes exist, for
example the extended ternary Golay code with generator matrix $[I_6\mid A]$, where
the rows of $A$ are $011111$, $101221$, $110122$, $121012$, $122101$, $112210$. Its
weight distribution is $1,264,440,24$ in weights $0,6,9,12$, so its minimum
distance is $6=n-k$, one less than MDS. MDS codes themselves exist in the
admissible range $n\le k+2$ (the tetracode). Lean proves both facts.
\end{remark}

\section{Lean formalization}
The directory \texttt{lean4/} is a self-contained Lean~4.19.0 project (core
library only, no Mathlib). $\F$ is \texttt{Fin 3}.
\begin{itemize}
\item \texttt{comb k c G j} $=\sum_{i<k}c_iG_{ij}$, \texttt{codeword k G m},
\texttt{wt n w} (Hamming weight on coordinates $<n$), \texttt{Nonzero k m}.
\texttt{IsGenMatrix n k G}: the rows are linearly independent (an $[n,k]$ code).
\texttt{InCode}, \texttt{MinDistGe}, \texttt{HasMinDist}, and
\texttt{IsMDS n k G} $:=$ \texttt{IsGenMatrix} $\wedge$ minimum distance
exactly $n-k+1$.
\item \texttt{dep} is Lemma~\ref{lem:dep}, proved by the elimination of its
proof (\texttt{elimV}, \texttt{skipCoord}). \texttt{cnt}, \texttt{pairWt\_eq}
and \texttt{sumTo\_bound} give Lemma~\ref{lem:count} and the exchange of sums.
\texttt{no\_mds\_ternary}: for $2\le k$, $k+3\le n$ and every $G$,
$\neg$\texttt{MinDistGe n k G (n-k+1)}. This is Theorem~\ref{thm:main}, proved
for all $n,k$ symbolically, not by enumeration.
\item \texttt{IsPAut}, \texttt{IsMAut}: permutations (lists of images) and
monomial maps sending every codeword to a codeword. \texttt{PAutOrder n k G N}
and \texttt{MAutOrder n k G N} say that a duplicate-free list of length $N$
consists of exactly the automorphisms.
\item \texttt{Conjecture1090} (and \texttt{Conjecture1090Mon}) is the conjunction
(attained $660$) $\wedge$ (all orders $\le660$) $\wedge$ (other orders divide
$132$). \texttt{conjecture\_00000001090\_false : $\neg$Conjecture1090} and
\texttt{conjecture\_00000001090\_monomial\_false} refute them through
\texttt{attained\_false}. \texttt{no\_mds\_12\_6}, \texttt{no\_dist7\_12\_6},
\texttt{no\_mds\_12\_6\_with} (any property $P$) and \texttt{no\_mds\_12\_7}
cover the readings above.
\item Non-vacuity: \texttt{tetra\_mds : IsMDS 4 2 tetra} and
\texttt{golay\_12\_6\_6} (the Golay code is a $[12,6]$ code with minimum distance
exactly $6$; checked over all $729$ messages by \texttt{decide +kernel}).
\item ``MDS'' dropped: \texttt{pair\_gen}, \texttt{pair\_aut},
\texttt{pairAuts\_nodup}, \texttt{atMost660\_false} (some $[12,6]$ code has
more than $660$ distinct permutation automorphisms), and
\texttt{pair\_order\_ge} (if \texttt{PAutOrder 12 6 pairG N} then
$N\ge768$). The exact value $46080$ is checked in \texttt{verify.py}
only.
\end{itemize}
\texttt{lake build} succeeds with no errors or warnings. There is no
\texttt{sorry}, no \texttt{native\_decide} and no added axiom.
\texttt{\#print axioms} shows only \texttt{propext} and \texttt{Quot.sound} for
all results about MDS codes. The two results on the pair code
(\texttt{atMost660\_false}, \texttt{pair\_order\_ge}) also use the standard axiom
\texttt{Classical.choice}, through core list lemmas.

\section{Independent check}
\texttt{verify.py} (Python~3 standard library) uses different methods:
\begin{itemize}
\item It searches all $1280^2$ pairs of words of weight $\ge7$ in $\F^8$ and finds
none spanning a code with all nonzero weights $\ge7$, so no $[8,2,7]$ code
exists. The best ternary $[8,2]$ code has $d=6$.
\item It checks that no systematic $[5,2,4]$ or $[6,3,4]$ ternary code exists.
An MDS code is always systematic on its first $k$ coordinates.
\item It evaluates the Griesmer bound ($n\ge14$ for $[n,6,7]_3$).
\item On $300$ random $[12,6]$ ternary codes, it runs the construction of
Theorem~\ref{thm:main} and computes the exact minimum distance (always
$\le6$).
\item It computes the weight distributions of the tetracode and the Golay code.
\item It enumerates all $46080$ pair-preserving permutations, confirms that they
are automorphisms of $P$, confirms that the weight-$2$ supports are the six pairs,
and re-checks the $768$ permutations used in Lean.
\end{itemize}

\section{Reproduction}
\begin{verbatim}
cd lean4 && lake build     # Lean 4.19.0, core only
cd .. && python3 verify.py
pdflatex report.tex && pdflatex report.tex
\end{verbatim}

\begin{thebibliography}{9}
\bibitem{ms} F.~J.~MacWilliams and N.~J.~A.~Sloane, \emph{The Theory of
Error-Correcting Codes}, North-Holland, 1977, Ch.~11 (MDS codes; $n\le q+k-1$
for $k\ge2$).
\bibitem{cs} J.~H.~Conway and N.~J.~A.~Sloane, \emph{Sphere Packings, Lattices
and Groups}, 3rd ed., Springer, 1999 (the ternary Golay code and $2.M_{12}$).
\end{thebibliography}
\end{document}
